\SourceLeaf{092}
ausgereicht hat; ich habe dazu das System der Begrenzungslinien der Flächentheile
mit positivem $T$ gewählt, es hätte aber eben so gut zu demselben Zweck die Betrachtung der Begrenzungslinien der Flächen mit positivem (oder negativem) $U$
dienen können.

\begin{center}\rule{0.16\linewidth}{0.4pt}\end{center}
\Needspace{8\baselineskip}
\begin{center}ZWEITE ABTHEILUNG.\end{center}

\FifthArticle{11}
Zur numerischen Bestimmung der Wurzeln solcher algebraischen Gleichungen, die nur aus drei Gliedern bestehen, lassen sich verschiedene Methoden anwenden, die hier einer Eleganz und Bequemlichkeit fähig werden, gegen welche
die mühsamen bei Gleichungen von weniger einfacher Gestalt unvermeidlichen
Operationen weit zurückstehen. Solche Methoden verdienen also wohl eine besondere Darstellung, zumal da Gleichungen von jener Form häufig genug vorkommen.

Es gilt dies zunächst von der Entwicklung der Wurzeln in unendliche Reihen. In der That lässt sich \textit{jede}, gleichviel ob reelle oder imaginäre, Wurzel einer Gleichung mit drei Gliedern durch eine convergente Reihe von einfachem
Fortschreitungsgesetz ausdrücken. Ich werde jedoch \textit{diese} Auflösungsart aus mehrern Gründen von meiner gegenwärtigen Betrachtung ganz ausschliessen, und bemerke hier nur, dass der Grad der Convergenz von dem gegenseitigen Verhalten
der Coëfficienten abhängig, dass sie desto langsamer ist, je näher dies Verhalten
demjenigen kommt, bei welchem die Gleichung zwei gleiche Wurzeln hat, und
dass in diesem Grenzfalle selbst sie schwächer ist, als bei irgendwelcher fallenden
geometrischen Progression. So bemerkenswerth auch diese Reihen in allgemeiner
theoretischer Rücksicht sind, so wird man doch, abgesehen von dem Falle, wo
ihre Convergenz eine sehr schnelle wird, in praktischer Beziehung immer den \textit{indirecten} Methoden den Vorzug geben, welche in den nachfolgenden Artikeln entwickelt werden sollen.

\FifthArticle{12}
Zur Auffindung der \textit{reellen} Wurzeln benutze ich meine im Jahre 1810 zuerst
gedruckte Hülfstafel für Logarithmen von Summen und Differenzen, oder, wo
eine grössere Genauigkeit verlangt wird, als Logarithmen mit fünf Zifern geben
\SourceLeaf{093}
können, die ähnliche aber erweiterte Tafel von \textsc{Matthiessen}. Ich habe ein paar
specielle Anwendungen dieses Verfahrens schon früher bekannt gemacht, nemlich
zur Auflösung der quadratischen Gleichungen bei der 1840 erschienenen zwanzigsten Ausgabe von \textsc{Vega}'s logarithmischem Handbuch, und zur Auflösung der
cubischen Gleichung; welche bei der parabolischen Bewegung zur Bestimmung
der wahren Anomalie dient, in Nro. 474 der Astronomischen Nachrichten. An
letzterm Orte ist auch bereits die allgemeine Anwendbarkeit des Verfahrens auf
alle algebraischen Gleichungen mit drei Gliedern bemerklich gemacht. Obgleich
nun die Ausführung dieses ganz elementarischen Gegenstandes gar keine Schwierigkeiten hat, so wird man doch, bei der ziemlich grossen Mannigfaltigkeit der
Fälle, einer übersichtlichen Sonderung derselben, und der Zusammenstellung der
gebrauchfertigen Vorschriften ein paar Seiten gern eingeräumt sehen.

Anstatt jener logarithmischen Hülfstafeln kann man sich auch der gewöhnlichen logarithmisch-trigonometrischen Tafeln bedienen: allein theils sind jene
im Allgemeinen für den gegenwärtigen Zweck von bequemerm Gebrauch, theils
gewähren sie doppelt so grosse Genauigkeit als die letztern. Ich würde daher die
Benutzung der trigonometrischen Tafeln für das in Rede stehende Geschäft auf
den seltenen Fall beschränken, wo man die durch siebenzifrige Logarithmen erreichbare Genauigkeit noch zu überschreiten wünscht und dazu die bekannten
zehnzifrigen Logarithmen in \textsc{Vlacq}'s oder \textsc{Vega}'s Thesaurus verwenden kann.
Uebrigens sind, wenn man sich der Hülfslogarithmen bedient, doppelt so viele
Fälle zu unterscheiden, als wenn die trigonometrischen Logarithmen gebraucht
werden. Als ein Nachtheil darf dies jedoch nicht angesehen werden: denn wenn
einmal die vollständige allgemeine Classification vorliegt, ist es leicht, jedem concreten Falle sein Fach anzuweisen, und das eigentliche indirecte Geschäft ist so
viel leichter auszuführen, wenn das ganze Fach nur den halben Umfang hat. Aber
gerade aus jenem Grunde ist für die Auflösung durch trigonometrische Logarithmen die allgemeine Classification kürzer und bequemer darzustellen, und ich werde
sie daher vorausschicken, da sodann die Classification für die andere Auflösungsform sich daraus von selbst ergibt.

\FifthArticle{13}
Die Ausführung der Methode wird, unmittelbar, nur auf Bestimmung der
\textit{positiven} Wurzeln einer vorgegebenen Gleichung gerichtet; die negativen ergeben
\SourceLeaf{094}
sich, indem man dasselbe Verfahren auf diejenige Gleichung anwendet, welche
aus jener durch Einführung der der ursprünglichen Unbekannten entgegengesetzten Grösse entsteht.

Die Gleichung setze ich in die Form
\[
x^{m+n}\pm e x^m\pm f=0
\]
wo $m$, $n$, $e$, $f$ gegebene \textit{positive} Grössen bedeuten. Diese Form umfasst eigentlich, nach Verschiedenheit der Combination der Zeichen, vier verschiedene Fälle,
wovon aber der erste, wo beidemal die oberen Zeichen gewählt werden, ausfällt,
da offenbar die Gleichung
\[
x^{m+n}+e x^m+f=0
\]
keine positive Wurzel haben kann. Uebrigens ist verstattet, vorauszusetzen, dass
$m$ und $n$ (worunter ganze Zahlen verstanden werden, obwohl die Anwendbarkeit
der Methode an sich davon unabhängig ist) keinen gemeinschaftlichen Divisor haben, indem auf diesen Fall jeder andere leicht zurückzuführen ist. Endlich werde
ich zur Abkürzung schreiben
\[
\frac{f^n}{e^{m+n}}=\lambda
\]

\Needspace{6\baselineskip}\begin{center}\textit{Erste Form.}\end{center}
\[
x^{m+n}+e x^m-f=0
\]
Indem man einen immer im ersten Quadranten zu nehmenden Winkel $\theta$
einführt, so dass
\[
\frac{x^{m+n}}{f}=\sin\theta^2,\quad
\frac{e x^m}{f}=\cos\theta^2
\]
wird, also (I)
\[
x^{m+n}=f\sin\theta^2,\quad
x^m=\frac{f\cos\theta^2}{e},\quad
x^n=e\mathrm{tang}\theta^2
\]
findet sich durch Elimination von $x$ die Gleichung
\[
\lambda=\frac{\sin\theta^{2m}}{\cos\theta^{2m+2n}}
\]
aus welcher $\theta$ bestimmt werden muss. Man erkennt leicht, dass der zweite Theil
dieser Gleichung als Function einer unbestimmten Grösse $\theta$ betrachtet, von $0$
\SourceLeaf{095}
bis $\infty$ wächst, während $\theta$ alle Werthe von $0$ bis $90^0$ durchläuft, und dass es
also einen, und nur einen Werth von $\theta$ gibt, der jener Gleichung Genüge leistet.
Nachdem derselbe gefunden ist, erhält man $x$ aus einer der Formeln I. Man bemerke, dass $\theta=45^0$ wird für $\lambda=2^n$, und dass folglich $\theta$ im ersten Octanten
zu suchen ist wenn $\lambda$ kleiner, im zweiten wenn $\lambda$ grösser ist als $2^n$.

\Needspace{6\baselineskip}\begin{center}\textit{Zweite Form.}\end{center}
\[
x^{m+n}-ex^m-f=0
\]
Man wird hier setzen
\[
fx^{-m-n}=\sin\theta^2,\quad ex^{-n}=\cos\theta^2
\]
oder (I)
\[
x^{m+n}=\frac{f}{\sin\theta^2},\quad
x^n=\frac{e}{\cos\theta^2},\quad
x^m=\frac{f\mathrm{cotang}\theta^2}{e}
\]
wonach also $\theta$ aus der Gleichung
\[
\lambda=\frac{\sin\theta^{2n}}{\cos\theta^{2m+2n}}
\]
zu bestimmen sein wird, was auf eine und nur auf eine Art geschehen kann: der
Werth von $x$ findet sich sodann durch eine der Gleichungen I. Im ersten oder
zweiten Octanten liegt $\theta$, jenachdem $\lambda$ kleiner oder grösser ist als $2^m$.

\Needspace{6\baselineskip}\begin{center}\textit{Dritte Form.}\end{center}
\[
x^{m+n}-ex^m+f=0
\]
Hier wird man setzen
\[
\frac{x^n}{e}=\sin\theta^2,\quad
\frac{fx^{-m}}{e}=\cos\theta^2
\]
oder (I)
\[
x^{m+n}=f\mathrm{tang}\theta^2,\quad
x^m=\frac{f}{e\cos\theta^2},\quad
x^n=e\sin\theta^2
\]
von welchen Formeln eine zur Bestimmung von $x$ dienen wird, sobald der Werth
von $\theta$ gefunden ist. Dieser ergibt sich durch Auflösung der Gleichung
\[
\lambda=\cos\theta^{2n}\sin\theta^{2m}
\]
Da das auf der rechten Seite stehende Glied dieser Gleichung, als Function einer
unbestimmten Grösse $\theta$ betrachtet, sowohl für $\theta=0$ als für $\theta=90^0$ verschwindet, so muss dazwischen ein grösster Werth liegen, und da das Differential des
\SourceLeaf{096}
Logarithmen dieser Function $=(2m\mathrm{cotg}\theta-2n\mathrm{tang}\theta)\,\mathrm{d}\theta$ ist, so findet der grösste
Werth Statt für $\theta=\theta^*$, wenn man $\surd\frac{m}{n}=\mathrm{tang}\theta^*$ setzt. Es wird demnach jene
Function von $0$ bis zu ihrem grössten Werthe, welcher offenbar
\[
=\frac{m^m n^n}{(m+n)^{m+n}}
\]
ist, zunehmen, und von da bis $0$ abnehmen, während $\theta$ von $0$ zu $\theta^*$ und von
da bis $90^0$ zunimmt. Der Maximumwerth ist daher jedenfalls grösser als der
Werth für $\theta=45^0$, d. i. grösser als $\frac{1}{2^{m+n}}$, den Fall ausgenommen, wo $m=n$,
und also $\frac{1}{2^{m+n}}$ selbst der Maximumwerth ist.

Man schliesst hieraus, dass jenachdem $\lambda$ grösser ist als
\[
\frac{m^m n^n}{(m+n)^{m+n}}
\]
oder kleiner, der Gleichung $\lambda=\cos\theta^{2n}\sin\theta^{2m}$ gar nicht oder durch zwei verschiedene Werthe von $\theta$ wird Genüge geleistet werden können. Im erstern Falle
hat die Gleichung $x^{m+n}-ex^m+f=0$ gar keine (positive) Wurzel, im andern
zwei. In dem speciellen Falle, wo
\[
\lambda=\frac{m^m n^n}{(m+n)^{m+n}}
\]
ist, fallen beide Auflösungen zusammen, und die Gleichung hat zwei gleiche Wurzeln, wofür man nach Gefallen eine der drei Formeln benutzen kann
\[
x^{m+n}=\frac{fm}{n},\quad
x^m=\frac{f(m+n)}{en},\quad
x^n=\frac{em}{m+n}
\]

Was übrigens in dem Falle, wo zwei Auflösungen wirklich vorhanden sind,
die Octanten betrifft, in welche die Werthe von $\theta$ fallen, so sieht man leicht,
dass wenn $\lambda$ grösser ist als $\frac{1}{2^{m+n}}$, beide Werthe von $\theta$ mit $\theta^*$ in demselben Octanten liegen, nemlich im ersten oder zweiten, jenachdem $m$ kleiner oder grösser
ist als $n$: ist hingegen $\lambda$ kleiner als $\frac{1}{2^{m+n}}$, so wird der eine Werth von $\theta$ im ersten, der andere im zweiten Octanten zu suchen sein. In dem speciellen Falle,
wo $\lambda=\frac{1}{2^{m+n}}$, ist $45^0$ selbst der eine Werth von $\theta$, und der andere liegt in demselben Octanten wie $\theta^*$.

Es mag noch die aus dieser Zergliederung aller drei Formen sich leicht ergebende Folge bemerkt werden, dass unsere Gleichung (insofern wir annehmen, dass
$m$ und $n$ keinen gemeinschaftlichen Divisor haben) nicht mehr als drei reelle
Wurzeln haben kann, was auch aus andern Gründen bekannt ist.
\SourceLeaf{097}
\FifthArticle{14}
Die vorstehenden Vorschriften werden nun leicht in diejenigen umgeschmolzen, die der Anwendung der Hülfslogarithmen entsprechen, da diese $A=\log a$.
$B=\log b$, $C=\log c$, betrachtet werden können wie die Logarithmen der Quadrate der Tangenten, Cosecanten und Secanten der von $45^0$ bis $90^0$ zunehmenden, oder, was dasselbe ist, wie die Logarithmen der Quadrate der Cotangenten,
Secanten und Cosecanten der von $45^0$ bis $0$ abnehmenden Winkel, also
\[
a=\mathrm{tang}\theta^2,\quad \frac{1}{b}=\sin\theta^2,\quad
\frac{1}{c}=\cos\theta^2
\]
für die Werthe von $\theta$ im zweiten Octanten, oder
\[
\frac{1}{a}=\mathrm{tang}\theta^2,\quad
\frac{1}{c}=\sin\theta^2,\quad
\frac{1}{b}=\cos\theta^2
\]
für die Werthe von $\theta$ im ersten Octanten.

Die vollständigen Vorschriften vereinige ich in folgendem Schema, wo eben
so wie oben
\[
\lambda=\frac{f^n}{e^{m+n}}
\]
gesetzt ist.

\Needspace{6\baselineskip}\begin{center}\textit{Erste Form.}\end{center}
\[
x^{m+n}+ex^m-f=0
\]
\textit{Erster Fall.} $\lambda>2^n$
\[
\lambda=a^{m+n}b^n=a^m c^n=\frac{c^{m+n}}{b^m}
\]
\[
x^{m+n}=\frac{f}{b},\quad x^m=\frac{f}{ec},\quad x^n=ea
\]
\textit{Zweiter Fall.} $\lambda<2^n$
\[
\lambda=\frac{b^n}{a^m}=\frac{c^n}{a^{m+n}}=\frac{b^{m+n}}{c^m}
\]
\[
x^{m+n}=\frac{f}{c},\quad x^m=\frac{f}{eb},\quad x^n=\frac{e}{a}
\]

\Needspace{6\baselineskip}\begin{center}\textit{Zweite Form.}\end{center}
\[
x^{m+n}-ex^m-f=0
\]
\textit{Erster Fall.} $\lambda>2^m$
\SourceLeaf{098}
\[
\lambda=a^{m+n}b^m=a^nc^m=\frac{c^{m+n}}{b^n}
\]
\[
x^{m+n}=fb,\quad x^m=\frac{f}{ea},\quad x^n=ec
\]
\textit{Zweiter Fall.}\quad $\lambda<2^m$
\[
\lambda=\frac{b^m}{a^n}=\frac{c^m}{a^{m+n}}=\frac{b^{m+n}}{c^n}
\]
\[
x^{m+n}=fc,\quad x^m=\frac{fa}{e},\quad x^n=eb
\]
\Needspace{6\baselineskip}\begin{center}\textit{Dritte Form.}\end{center}
\[
x^{m+n}-ex^m+f=0
\]
\textit{Erster Fall.}\quad $\frac{1}{\lambda}<\frac{(m+n)^{m+n}}{m^mn^n}$

Gar keine Auflösung.

\textit{Zweiter Fall.}\quad $\frac{1}{\lambda}=\frac{(m+n)^{m+n}}{m^mn^n}$

Zwei gleiche Wurzeln, zu deren Bestimmung eine der Gleichungen
\[
x^{m+n}=\frac{fm}{n},\quad x^m=\frac{f(m+n)}{en},\quad x^n=\frac{em}{m+n}
\]
dient.

\textit{Dritter Fall.}\quad $\frac{1}{\lambda}$ grösser als $\frac{(m+n)^{m+n}}{m^mn^n}$ aber nicht grösser als $2^{m+n}$, und zugleich
$m$ grösser als $n$.

Zwei Wurzeln, für welche
\[
\frac{1}{\lambda}=a^nb^{m+n}=\frac{c^{m+n}}{a^m}=b^mc^n
\]
\[
x^{m+n}=fa,\quad x^m=\frac{fc}{e},\quad x^n=\frac{e}{b}
\]
\textit{Vierter Fall.}\quad Für $\frac{1}{\lambda}$ dieselben Grenzen, wie im dritten Fall, aber $m$ kleiner
als $n$.

Zwei Wurzeln, für welche
\[
\frac{1}{\lambda}=a^mb^{m+n}=\frac{c^{m+n}}{a^n}=b^nc^m
\]
\[
x^{m+n}=\frac{f}{a},\quad x^m=\frac{fb}{e},\quad x^n=\frac{e}{c}
\]
\textit{Fünfter Fall.}\quad $\frac{1}{\lambda}$ grösser als $2^{m+n}$.

Zwei Wurzeln, wovon die eine durch die Formeln des dritten Falles, die
andere durch die des vierten bestimmt wird.
\SourceLeaf{099}
Es mag noch bemerkt werden, dass im dritten Falle der Werth von $a$, welcher der einen Wurzel entspricht, kleiner als $\frac{m}{n}$, der zur andern Wurzel gehörende grösser als $\frac{m}{n}$ ist; im vierten Falle verhalten sich die beiden Werthe von
$a$ auf ähnliche Weise gegen $\frac{n}{m}$.

\FifthArticle{15}
Ueber die Anwendung dieser Vorschriften ist noch folgendes beizufügen.

Zur Bestimmung jeder Wurzel sind zwei Operationen auszuführen: zuerst
aus $\lambda$ den dazu gehörenden Werth von $a$ (und damit zugleich den von $b$ oder $c$)
abzuleiten; sodann, aus diesem,den Werth von $x$ zu berechnen. Für jede dieser beiden Operationen kann man unter drei Formeln wählen; ich ziehe in den
meisten Fällen die zuerst angesetzten vor. Bei allen diesen Rechnungen hat man
es gar nicht mit den Grössen $\lambda$, $a$, $b$, $c$ selbst, sondern nur mit ihren Logarithmen zu thun. Die erste Operation ist eine indirecte, und beruhet demnach in der
Regel auf mehrern stufenweise fortschreitenden Annäherungen, wobei es bequem
gefunden werden wird, zu Anfang Tafeln mit einer geringern Anzahl von Zifern
zu gebrauchen. \textsc{Matthiessens} Tafel hat bekanntlich sieben Decimalen: die meinige fünf; \textsc{Encke} und \textsc{Ursin} haben sie mit vier Zifern abdrucken lassen, und wenn
man beim Anfange der Arbeit noch gar keine Kenntniss einer ersten groben Annäherung mitbringt, wird man es vielleicht vortheilhaft finden, einen noch kürzern Extract der Tafeln mit nur drei Zifern auf einem besondern Blättchen vor
sich zu haben, etwa so:
\[
\begin{array}[t]{c|c|}
A&B\\\hline
0&0{,}301\\
0{,}1&0{,}254\\
0{,}2&0{,}212\\
0{,}3&0{,}176\\
0{,}4&0{,}146\\
0{,}5&0{,}119\\
0{,}6&0{,}097\\
0{,}7&0{,}079\\
0{,}8&0{,}064\\
0{,}9&0{,}051\\
1{,}0&0{,}041
\end{array}
\begin{array}[t]{c|c|}
A&B\\\hline
1{,}0&0{,}041\\
1{,}1&0{,}033\\
1{,}2&0{,}027\\
1{,}3&0{,}021\\
1{,}4&0{,}017\\
1{,}5&0{,}014\\
1{,}6&0{,}011\\
1{,}7&0{,}009\\
1{,}8&0{,}007\\
1{,}9&0{,}005\\
2{,}0&0{,}004
\end{array}
\begin{array}[t]{c|c}
A&B\\\hline
2{,}0&0{,}004\\
2{,}1&0{,}003\\
2{,}2&0{,}003\\
2{,}3&0{,}002\\
2{,}4&0{,}002\\
2{,}5&0{,}001\\
2{,}9&0{,}001\\
3{,}0&0{,}000
\end{array}
\]
\SourceLeaf{100}
\FifthArticle{16}
Als Beispiel mag die Gleichung
\[
x^7+28x^4-480=0
\]
dienen, wo
\[
\lambda=\frac{6750}{823543},\quad \log\frac{1}{\lambda}=2{,}0863825
\]
wird. Die Gleichung hat die erste Form, mithin \textit{eine} positive Wurzel, und gehört, da $\lambda$ kleiner ist als $8$, zum zweiten Fall. Die erste Operation besteht darin,
dass der Gleichung $\log\frac{1}{\lambda}=4A-3B$ Genüge geschehe, also, wenn man die
Rechnung mit drei Decimalen anfängt, dieser
\[
2{,}086=4A-3B
\]
Ein flüchtiger Blick auf obige Tafel zeigt schon, dass $A$ zwischen $0{,}5$ und $0{,}6$
zu suchen sei. Es wird nämlich
\[
\begin{array}{c|c|c}
A&4A-3B&\mbox{Fehler}\\\hline
0{,}5&1{,}643&-0{,}443\\
0{,}6&2{,}109&+0{,}023
\end{array}
\]
woraus sich auf einen genauern Werth $0{,}595$ schliessen lässt. Eine neue Rechnung nach den Tafeln mit fünf Decimalen, wo also $\log\frac{1}{\lambda}=2{,}08638$ zu setzen
ist, gibt
\[
\begin{array}{c|c|c}
A&4A-3B&\mbox{Fehler}\\\hline
0{,}595&2{,}08501&-0{,}00137\\
0{,}596&2{,}08961&+0{,}00323
\end{array}
\]
woraus der noch genauere Werth $0{,}5953$ erkannt wird. Endlich für sieben Decimalen hat man
\[
\begin{array}{c|c|c}
A&4A-3B&\mbox{Fehler}\\\hline
0{,}5952&2{,}0859279&-0{,}0004546\\
0{,}5953&2{,}0863885&+0{,}0000060
\end{array}
\]
Zu dem Werthe $A=0{,}5953$ muss also noch die Correction $-\frac{60}{4606}$ Einheiten
der vierten Decimale hinzukommen, in welcher Form ich sie beibehalte, da es,
wenn zur Bestimmung von $x$ die erste Formel
\SourceLeaf{101}
\[
x^{7}=\frac{f}{c}
\]
gebraucht werden soll, nur darauf ankommt, den entsprechenden Werth von $C$
zu finden. Diesen erhält man, indem man zu dem neben $A=0{,}5953$ stehenden Werthe $C=0{,}6935705$ die Correction $-\frac{60}{4606}\times798$ hinzufügt, letztere
wie Einheiten der siebenten Decimale betrachtet, also
\[
\begin{array}{rcl}
C&=&0{,}6935695\\
\log f&=&2{,}6812412\\ \hline
7\log x&=&1{,}9876717\\
\log x&=&0{,}2839531\\
x&=&1{,}9228841
\end{array}
\]

Zur Auffindung der negativen Wurzeln wird man $x=-y$ schreiben und
die positiven Wurzeln der Gleichung
\[
y^{7}-28y^{4}+480=0
\]
aufsuchen. Diese gehört zur dritten Form, und da $\frac{1}{\lambda}=\frac{823543}{6750}$ grösser ist als
$\frac{7^{7}}{3^{3}4^{4}}=\frac{823543}{6912}$, aber kleiner als $2^{7}=128$, zugleich auch $m$ grösser ist als $n$, so
gilt der dritte Fall, oder es finden zwei Wurzeln Statt, zu deren Ausmittlung
der Gleichung
\[
2{,}0863825=3A+7B
\]
genügt werden muss. Aus der Schlussbemerkung des 14. Art. weiss man, dass
der eine Werth von $A$ kleiner, der andere grösser sein muss als $\log\frac{4}{3}=0{,}12494$.
Auch ergeben sich die Grenzen der Werthe von $A$ sofort aus der obigen Tafel
mit dreizifrigen Logarithmen, nach welchen man erhält:
\[
\begin{array}{c|c|r}
A&3A+7B&\mbox{Fehler}\\ \hline
0{,}0&2{,}107&+0{,}021\\
0{,}1&2{,}078&-0{,}008\\
0{,}2&2{,}084&-0{,}002\\
0{,}3&2{,}132&+0{,}046
\end{array}
\]
Will man zur nähern Bestimmung zuerst vierzifrige Logarithmen gebrauchen, so
hat man zunächst für die erste Auflösung
